% Options for packages loaded elsewhere
% Options for packages loaded elsewhere
\PassOptionsToPackage{unicode}{hyperref}
\PassOptionsToPackage{hyphens}{url}
\PassOptionsToPackage{dvipsnames,svgnames,x11names}{xcolor}
%
\documentclass[
  spanish,
  letterpaper,
  DIV=11,
  numbers=noendperiod]{scrreprt}
\usepackage{xcolor}
\usepackage{amsmath,amssymb}
\setcounter{secnumdepth}{0}
\usepackage{iftex}
\ifPDFTeX
  \usepackage[T1]{fontenc}
  \usepackage[utf8]{inputenc}
  \usepackage{textcomp} % provide euro and other symbols
\else % if luatex or xetex
  \usepackage{unicode-math} % this also loads fontspec
  \defaultfontfeatures{Scale=MatchLowercase}
  \defaultfontfeatures[\rmfamily]{Ligatures=TeX,Scale=1}
\fi
\usepackage{lmodern}
\ifPDFTeX\else
  % xetex/luatex font selection
  \setmonofont[]{JuliaMono}
\fi
% Use upquote if available, for straight quotes in verbatim environments
\IfFileExists{upquote.sty}{\usepackage{upquote}}{}
\IfFileExists{microtype.sty}{% use microtype if available
  \usepackage[]{microtype}
  \UseMicrotypeSet[protrusion]{basicmath} % disable protrusion for tt fonts
}{}
\makeatletter
\@ifundefined{KOMAClassName}{% if non-KOMA class
  \IfFileExists{parskip.sty}{%
    \usepackage{parskip}
  }{% else
    \setlength{\parindent}{0pt}
    \setlength{\parskip}{6pt plus 2pt minus 1pt}}
}{% if KOMA class
  \KOMAoptions{parskip=half}}
\makeatother
% Make \paragraph and \subparagraph free-standing
\makeatletter
\ifx\paragraph\undefined\else
  \let\oldparagraph\paragraph
  \renewcommand{\paragraph}{
    \@ifstar
      \xxxParagraphStar
      \xxxParagraphNoStar
  }
  \newcommand{\xxxParagraphStar}[1]{\oldparagraph*{#1}\mbox{}}
  \newcommand{\xxxParagraphNoStar}[1]{\oldparagraph{#1}\mbox{}}
\fi
\ifx\subparagraph\undefined\else
  \let\oldsubparagraph\subparagraph
  \renewcommand{\subparagraph}{
    \@ifstar
      \xxxSubParagraphStar
      \xxxSubParagraphNoStar
  }
  \newcommand{\xxxSubParagraphStar}[1]{\oldsubparagraph*{#1}\mbox{}}
  \newcommand{\xxxSubParagraphNoStar}[1]{\oldsubparagraph{#1}\mbox{}}
\fi
\makeatother

\usepackage{color}
\usepackage{fancyvrb}
\newcommand{\VerbBar}{|}
\newcommand{\VERB}{\Verb[commandchars=\\\{\}]}
\DefineVerbatimEnvironment{Highlighting}{Verbatim}{commandchars=\\\{\}}
% Add ',fontsize=\small' for more characters per line
\usepackage{framed}
\definecolor{shadecolor}{RGB}{241,243,245}
\newenvironment{Shaded}{\begin{snugshade}}{\end{snugshade}}
\newcommand{\AlertTok}[1]{\textcolor[rgb]{0.68,0.00,0.00}{#1}}
\newcommand{\AnnotationTok}[1]{\textcolor[rgb]{0.37,0.37,0.37}{#1}}
\newcommand{\AttributeTok}[1]{\textcolor[rgb]{0.40,0.45,0.13}{#1}}
\newcommand{\BaseNTok}[1]{\textcolor[rgb]{0.68,0.00,0.00}{#1}}
\newcommand{\BuiltInTok}[1]{\textcolor[rgb]{0.00,0.23,0.31}{#1}}
\newcommand{\CharTok}[1]{\textcolor[rgb]{0.13,0.47,0.30}{#1}}
\newcommand{\CommentTok}[1]{\textcolor[rgb]{0.37,0.37,0.37}{#1}}
\newcommand{\CommentVarTok}[1]{\textcolor[rgb]{0.37,0.37,0.37}{\textit{#1}}}
\newcommand{\ConstantTok}[1]{\textcolor[rgb]{0.56,0.35,0.01}{#1}}
\newcommand{\ControlFlowTok}[1]{\textcolor[rgb]{0.00,0.23,0.31}{\textbf{#1}}}
\newcommand{\DataTypeTok}[1]{\textcolor[rgb]{0.68,0.00,0.00}{#1}}
\newcommand{\DecValTok}[1]{\textcolor[rgb]{0.68,0.00,0.00}{#1}}
\newcommand{\DocumentationTok}[1]{\textcolor[rgb]{0.37,0.37,0.37}{\textit{#1}}}
\newcommand{\ErrorTok}[1]{\textcolor[rgb]{0.68,0.00,0.00}{#1}}
\newcommand{\ExtensionTok}[1]{\textcolor[rgb]{0.00,0.23,0.31}{#1}}
\newcommand{\FloatTok}[1]{\textcolor[rgb]{0.68,0.00,0.00}{#1}}
\newcommand{\FunctionTok}[1]{\textcolor[rgb]{0.28,0.35,0.67}{#1}}
\newcommand{\ImportTok}[1]{\textcolor[rgb]{0.00,0.46,0.62}{#1}}
\newcommand{\InformationTok}[1]{\textcolor[rgb]{0.37,0.37,0.37}{#1}}
\newcommand{\KeywordTok}[1]{\textcolor[rgb]{0.00,0.23,0.31}{\textbf{#1}}}
\newcommand{\NormalTok}[1]{\textcolor[rgb]{0.00,0.23,0.31}{#1}}
\newcommand{\OperatorTok}[1]{\textcolor[rgb]{0.37,0.37,0.37}{#1}}
\newcommand{\OtherTok}[1]{\textcolor[rgb]{0.00,0.23,0.31}{#1}}
\newcommand{\PreprocessorTok}[1]{\textcolor[rgb]{0.68,0.00,0.00}{#1}}
\newcommand{\RegionMarkerTok}[1]{\textcolor[rgb]{0.00,0.23,0.31}{#1}}
\newcommand{\SpecialCharTok}[1]{\textcolor[rgb]{0.37,0.37,0.37}{#1}}
\newcommand{\SpecialStringTok}[1]{\textcolor[rgb]{0.13,0.47,0.30}{#1}}
\newcommand{\StringTok}[1]{\textcolor[rgb]{0.13,0.47,0.30}{#1}}
\newcommand{\VariableTok}[1]{\textcolor[rgb]{0.07,0.07,0.07}{#1}}
\newcommand{\VerbatimStringTok}[1]{\textcolor[rgb]{0.13,0.47,0.30}{#1}}
\newcommand{\WarningTok}[1]{\textcolor[rgb]{0.37,0.37,0.37}{\textit{#1}}}

\usepackage{longtable,booktabs,array}
\newcounter{none} % for unnumbered tables
\usepackage{calc} % for calculating minipage widths
% Correct order of tables after \paragraph or \subparagraph
\usepackage{etoolbox}
\makeatletter
\patchcmd\longtable{\par}{\if@noskipsec\mbox{}\fi\par}{}{}
\makeatother
% Allow footnotes in longtable head/foot
\IfFileExists{footnotehyper.sty}{\usepackage{footnotehyper}}{\usepackage{footnote}}
\makesavenoteenv{longtable}
\usepackage{graphicx}
\makeatletter
\newsavebox\pandoc@box
\newcommand*\pandocbounded[1]{% scales image to fit in text height/width
  \sbox\pandoc@box{#1}%
  \Gscale@div\@tempa{\textheight}{\dimexpr\ht\pandoc@box+\dp\pandoc@box\relax}%
  \Gscale@div\@tempb{\linewidth}{\wd\pandoc@box}%
  \ifdim\@tempb\p@<\@tempa\p@\let\@tempa\@tempb\fi% select the smaller of both
  \ifdim\@tempa\p@<\p@\scalebox{\@tempa}{\usebox\pandoc@box}%
  \else\usebox{\pandoc@box}%
  \fi%
}
% Set default figure placement to htbp
\def\fps@figure{htbp}
\makeatother



\ifLuaTeX
\usepackage[bidi=basic,shorthands=off]{babel}
\else
\usepackage[bidi=default,shorthands=off]{babel}
\fi
\ifLuaTeX
  \usepackage{selnolig} % disable illegal ligatures
\fi


\setlength{\emergencystretch}{3em} % prevent overfull lines

\providecommand{\tightlist}{%
  \setlength{\itemsep}{0pt}\setlength{\parskip}{0pt}}



 


\usepackage{mathrsfs}
\usepackage{amsthm}
\usepackage{mathpartir}
\usepackage{tikz}
\newcommand{\lean}[1]{\texttt{#1}}
%\usepackage{twemojis}
%\usepackage{myunixode}
\usepackage[normalem]{ulem}
\usepackage{scrlayer-scrpage}
\usepackage[most]{tcolorbox}
\automark{section}
\setkomafont{pagefoot}{\upshape}
\cfoot*{}
\ofoot*{\pagemark}

\makeatletter
\def\redsquiggly{\bgroup \markoverwith{\textcolor{red}{\lower3.5\p@\hbox{\sixly \char58}}}\ULon}
\def\brownsquiggly{\bgroup \markoverwith{\textcolor[HTML]{B8860B}{\lower3.5\p@\hbox{\sixly \char58}}}\ULon}
\def\bluesquiggly{\bgroup \markoverwith{\textcolor[HTML]{1E90FF}{\lower3.5\p@\hbox{\sixly \char58}}}\ULon}
\makeatother

\renewcommand{\NormalTok}[1]{\textcolor[HTML]{000000}{#1}}
\renewcommand{\KeywordTok}[1]{\textcolor[HTML]{0000FF}{#1}}
\renewcommand{\SpecialCharTok}[1]{}
\renewcommand{\ErrorTok}[1]{\redsquiggly{#1}}
\renewcommand{\WarningTok}[1]{\redsquiggly{\textcolor[HTML]{0000FF}{#1}}}
\renewcommand{\StringTok}[1]{\textcolor[HTML]{A52A2A}{#1}}
\renewcommand{\CommentTok}[1]{\textcolor[HTML]{008000}{#1}}
\renewcommand{\InformationTok}[1]{\textcolor[HTML]{D2691E}{\textbf{#1}}}
\renewcommand{\RegionMarkerTok}[1]{▼\:\textcolor[HTML]{008000}{\textbf{#1}}}
\renewcommand{\SpecialStringTok}[1]{\textcolor[HTML]{4682B4}{\textbf{#1}}}
\renewcommand{\ConstantTok}[1]{\textcolor[HTML]{DC143C}{#1}}
\renewcommand{\AnnotationTok}[1]{\brownsquiggly{#1}}
\renewcommand{\AlertTok}[1]{\brownsquiggly{\textcolor[HTML]{0000FF}{#1}}}
\renewcommand{\OtherTok}[1]{\bluesquiggly{#1}}
\renewcommand{\DocumentationTok}[1]{\bluesquiggly{\textcolor[HTML]{0000FF}{#1}}}

\newenvironment{ind}
	{\begin{list}{}{\setlength{\leftmargin}{1em}}\item\relax}
	{\end{list}}

%redefines Shaded so it can't break
\newcommand{\nobreakShaded}{\renewenvironment{Shaded}
	{\begin{tcolorbox}[frame hidden, boxrule=0pt, colback=shadecolor,
		sharp corners, left=0pt, right=0pt, top=0pt, bottom=0pt]}
	{\end{tcolorbox}}}

%Make end of environment ignore pars that come after it
\def\useignorespacesandallpars#1\ignorespaces\fi{%
#1\fi\ignorespacesandallpars}

\makeatletter
\def\ignorespacesandallpars{%
  \@ifnextchar\par
    {\expandafter\ignorespacesandallpars\@gobble}%
    {}%
}
\makeatother

\newenvironment{inpt}
	{\nobreakShaded\noindent\begin{minipage}[t]{0.64\textwidth}
		\uline{Código Lean}}
	{\end{minipage}\hfill\useignorespacesandallpars}

\newenvironment{outpt}
	{\nobreakShaded\begin{minipage}[t]{0.34\textwidth}
		\uline{Estado de la prueba}}
	{\end{minipage}\\}

\newenvironment{bef}
	{\nobreakShaded\noindent\begin{minipage}[t]{0.49\textwidth}
		\uline{Tactic State Before Using Strategy}}
	{\end{minipage}\hfill\useignorespacesandallpars}

\newenvironment{aft}
	{\nobreakShaded\begin{minipage}[t]{0.49\textwidth}
		\uline{Tactic State After Using Strategy}}
	{\end{minipage}\\}

\newenvironment{numex}[1]
	{\begin{minipage}[t]{0.04\textwidth}\vspace{3pt}{#1}.
		\end{minipage}\nobreakShaded\begin{minipage}[t]{0.96\textwidth}\vspace{0pt}}
	{\end{minipage}}

\newenvironment{mdsk}
	{\medskip}
	{}

\newenvironment{absnobreak}
  {\par\nobreak\vfil\penalty0\vfilneg
   \vtop\bgroup}
  {\par\xdef\tpd{\the\prevdepth}\egroup
   \prevdepth=\tpd}

\newcommand{\excl}[1]{}
\newcommand{\incl}[1]{#1}
\newcommand{\lrg}[1]{{\Large #1}}

\newcommand{\setmin}{\mathbin{\backslash}}
\newcommand{\symmdiff}{\mathbin{∆}}

\pagenumbering{roman}  %So front matter uses roman numerals.  Switch back to arabic at beginning of preface.


%Adds gray background to inline code, but inline code won't break across lines.
%\let\oldtexttt\texttt
%\renewcommand{\texttt}[1]{\colorbox{shadecolor}{\oldtexttt{#1}}}
\KOMAoption{captions}{tableheading}
\makeatletter
\@ifpackageloaded{bookmark}{}{\usepackage{bookmark}}
\makeatother
\makeatletter
\@ifpackageloaded{caption}{}{\usepackage{caption}}
\AtBeginDocument{%
\ifdefined\contentsname
  \renewcommand*\contentsname{Tabla de contenidos}
\else
  \newcommand\contentsname{Tabla de contenidos}
\fi
\ifdefined\listfigurename
  \renewcommand*\listfigurename{Listado de Figuras}
\else
  \newcommand\listfigurename{Listado de Figuras}
\fi
\ifdefined\listtablename
  \renewcommand*\listtablename{Listado de Tablas}
\else
  \newcommand\listtablename{Listado de Tablas}
\fi
\ifdefined\figurename
  \renewcommand*\figurename{Figura}
\else
  \newcommand\figurename{Figura}
\fi
\ifdefined\tablename
  \renewcommand*\tablename{Tabla}
\else
  \newcommand\tablename{Tabla}
\fi
}
\@ifpackageloaded{float}{}{\usepackage{float}}
\floatstyle{ruled}
\@ifundefined{c@chapter}{\newfloat{codelisting}{h}{lop}}{\newfloat{codelisting}{h}{lop}[chapter]}
\floatname{codelisting}{Listado}
\newcommand*\listoflistings{\listof{codelisting}{Listado de Listados}}
\makeatother
\makeatletter
\makeatother
\makeatletter
\@ifpackageloaded{caption}{}{\usepackage{caption}}
\@ifpackageloaded{subcaption}{}{\usepackage{subcaption}}
\makeatother
\usepackage{bookmark}
\IfFileExists{xurl.sty}{\usepackage{xurl}}{} % add URL line breaks if available
\urlstyle{same}
% fallback for those not using the hyperref driver hyperxmp:
\makeatletter
\@ifundefined{xmpquote}{\newcommand{\xmpquote}[1]{#1}}{}
\makeatother
\hypersetup{
  pdftitle={Una prueba transicional de la confluencia de la reducción \textbackslash eta},
  pdfauthor={Maximiliano Onofre Martínez},
  pdflang={es},
  colorlinks=true,
  linkcolor={blue},
  filecolor={Maroon},
  citecolor={Blue},
  urlcolor={Blue},
  pdfcreator={LaTeX via pandoc}}


\title{Una prueba transicional de la confluencia de la reducción
\(\eta\)}
\author{Maximiliano Onofre Martínez}
\date{}
\begin{document}
\maketitle

%Can't be in preamble because Quarto loads amsthm too late.
\theoremstyle{definition}
\newtheorem*{thm}{Theorem}
\newcommand{\thmnm}{Theorem}
\newtheorem*{namedthm}{\thmnm}
\newtheorem*{dfn}{Definition}
\newcommand{\defnm}{Definition}
\newtheorem*{nameddfn}{\defnm}

\newenvironment{nthm}[1]
  {\renewcommand{\thmnm}{#1}\begin{namedthm}}
  {\end{namedthm}}

\newenvironment{ndfn}[1]
  {\renewcommand{\defnm}{#1}\begin{nameddfn}}
  {\end{nameddfn}}

\newenvironment{npf}[1]
  {\begin{proof}[#1]}
  {\end{proof}}

% Usage:
% ::: {.thm} will create an unnumbered thm environment with title Theorem.
% ::: {.nthm arguments="Name of Theorem"} will create an unnumbered theorem whose title is Name of Theorem.

\let\oldgreater\textgreater
\renewcommand{\textgreater}{\null\oldgreater}   % To prevent => changing to double arrow

\renewcommand*\contentsname{Tabla de contenidos}
{
\setcounter{tocdepth}{1}
\tableofcontents
}

\bookmarksetup{startatroot}

\chapter*{Introducción}\label{introducciuxf3n}
\addcontentsline{toc}{chapter}{Introducción}

\markboth{Introducción}{Introducción}

\pagenumbering{arabic}

TODO

\bookmarksetup{startatroot}

\chapter{Relaciones}\label{relaciones}

En teoría de conjuntos se entiende una relación binaria \(R\) sobre un
conjunto \(A\) como un conjunto de pares ordenados, de modo que
\(R(a,b)\) significa \((a,b) \in R\). Una alternativa es pensar en \(R\)
como una función que, al aplicarse a \(a\) y \(b\), produce una
proposición que expresa que ambos elementos están relacionados.
Adoptaremos esta presentación en Lean mediante la siguiente definición:

\begin{Shaded}
\begin{Highlighting}[]
\KeywordTok{def}\NormalTok{ relation (X : }\KeywordTok{Type}\NormalTok{) := X → X → }\KeywordTok{Prop}
\end{Highlighting}
\end{Shaded}

Así, si \texttt{(R\ :\ relation\ X)} y \texttt{(a\ b\ :\ X)}, la
expresión \texttt{(R\ a\ b)} representa la proposición que afirma que
\texttt{a} está relacionado con \texttt{b} mediante \texttt{R}.

Sea \texttt{R\ :\ relation\ X}. La relación \texttt{R} es
\emph{reflexiva} si todo término de tipo \texttt{X} está relacionado
consigo mismo, y es \emph{transitiva} si, para cualesquiera términos
\texttt{a}, \texttt{b} y \texttt{c} de tipo \texttt{X}, de
\texttt{R\ a\ b} y \texttt{R\ b\ c} se sigue \texttt{R\ a\ c}. Estas
propiedades pueden expresarse en Lean de la siguiente manera:

\begin{Shaded}
\begin{Highlighting}[]
\NormalTok{variable \{X : }\KeywordTok{Type}\NormalTok{\}}

\KeywordTok{def}\NormalTok{ reflexive (R : relation X) : }\KeywordTok{Prop}\NormalTok{ :=}
\NormalTok{  ∀ a : X, R a a}

\KeywordTok{def}\NormalTok{ transitive (R : relation X) : }\KeywordTok{Prop}\NormalTok{ :=}
\NormalTok{  ∀ a b c : X, R a b → R b c → R a c}
\end{Highlighting}
\end{Shaded}

Las llaves en \texttt{variable\ \{X\ :\ Type\}} indican que \texttt{X}
es un argumento \emph{implícito}: en lugar de exigir que se proporcione,
Lean lo infiere a partir del contexto. Así, como \texttt{R} tiene tipo
\texttt{relation\ X}, Lean determina \texttt{X} a partir del tipo de
\texttt{R}, y basta escribir \texttt{reflexive\ R} o
\texttt{transitive\ R} sin indicar el tipo sobre el cual está definida
la relación.

\section{1.1. Clausuras}\label{clausuras}

A continuación presentamos las tres clausuras de \(R\) que utilizaremos,
junto con sus reglas de inferencia y sus definiciones en Lean.
Interpretaremos \(R(a,b)\) como un paso de \(a\) a \(b\).

I. La clausura reflexiva \(R^{=}\) relaciona dos elementos si son
iguales o si están relacionados por \(R\). Se define mediante las
siguientes reglas:

\[
\inferrule*[right=\textsc{refl}]
  { }
  {R^=(a,a)}
\qquad
\inferrule*[right=\textsc{step}]
  {R(a,b)}
  {R^=(a,b)}
\]

\begin{Shaded}
\begin{Highlighting}[]
\NormalTok{inductive clos\_refl (R : relation X) : relation X}
\NormalTok{  | refl : ∀ a : X,}
\NormalTok{      clos\_refl R a a}
\NormalTok{  | step : ∀ a b : X,}
\NormalTok{      R a b → clos\_refl R a b}
\end{Highlighting}
\end{Shaded}

II. La clausura transitiva \(R^{+}\) relaciona dos elementos cuando
pueden conectarse mediante una sucesión finita de uno o más pasos de
\(R\). Sus reglas de inferencia son:

\[
\inferrule*[right=\textsc{step}]
  {R(a,b)}
  {R^+(a,b)}
\qquad
\inferrule*[right=\textsc{trans}]
  {R^+(a,b) \\
   R^+(b,c)}
  {R^+(a,c)}
\]

\begin{Shaded}
\begin{Highlighting}[]
\NormalTok{inductive clos\_trans (R : relation X) : relation X}
\NormalTok{  | step : ∀ a b : X,}
\NormalTok{      R a b → clos\_trans R a b}
\NormalTok{  | trans : ∀ a b c : X,}
\NormalTok{      clos\_trans R a b → clos\_trans R b c →}
\NormalTok{      clos\_trans R a c}
\end{Highlighting}
\end{Shaded}

III. La clausura reflexiva-transitiva \(R^*\) relaciona dos elementos
cuando pueden conectarse mediante una sucesión finita de cero o más
pasos de \(R\). Una sucesión de cero pasos comienza y termina en el
mismo elemento.

\[
\inferrule*[right=\textsc{refl}]
  { }
  {R^*(a,a)}
\qquad
\inferrule*[right=\textsc{step}]
  {R(a,b)}
  {R^*(a,b)}
\qquad
\inferrule*[right=\textsc{trans}]
  {R^*(a,b) \\
   R^*(b,c)}
  {R^*(a,c)}
\]

\begin{Shaded}
\begin{Highlighting}[]
\NormalTok{inductive clos\_refl\_trans (R : relation X) : relation X}
\NormalTok{  | refl : ∀ a : X,}
\NormalTok{      clos\_refl\_trans R a a}
\NormalTok{  | step : ∀ a b : X,}
\NormalTok{      R a b → clos\_refl\_trans R a b}
\NormalTok{  | trans : ∀ a b c : X,}
\NormalTok{      clos\_refl\_trans R a b → clos\_refl\_trans R b c →}
\NormalTok{      clos\_refl\_trans R a c}
\end{Highlighting}
\end{Shaded}

\textbf{Lema 1.1.1.} La clausura reflexiva-transitiva de toda relación
es igual a la clausura transitiva de su clausura reflexiva.
\[ R^* = (R^{=})^+ \] Su enunciado en Lean es:

\begin{Shaded}
\begin{Highlighting}[]
\KeywordTok{lemma}\NormalTok{ clos\_refl\_trans\_eq : ∀ R : relation X, }
\NormalTok{    clos\_refl\_trans R = clos\_trans (clos\_refl R)}
\end{Highlighting}
\end{Shaded}

\section{1.2. Propiedad del diamante y
confluencia}\label{propiedad-del-diamante-y-confluencia}

Una relación \(R\) tiene la \emph{propiedad del diamante} si y solo si
\[
\forall m,m_1,m_2,\quad
R(m,m_1)\implies R(m,m_2)\implies
\exists m_3,\;
\bigl[R(m_1,m_3)\land R(m_2,m_3)\bigr].
\]

En Lean:

\begin{Shaded}
\begin{Highlighting}[]
\KeywordTok{def}\NormalTok{ diamond (R : relation X) : }\KeywordTok{Prop}\NormalTok{ :=}
\NormalTok{  ∀ m m₁ m₂ : X, R m m₁ → R m m₂ → ∃ m₃ : X, R m₁ m₃ ∧ R m₂ m₃}
\end{Highlighting}
\end{Shaded}

Escribiendo \(a \longrightarrow b\) en lugar de \(R(a,b)\), gráficamente
esto significa que siempre podemos completar el siguiente diagrama:

\begin{center}
\begin{tikzpicture}[>=stealth, scale=0.7]
  \node (m)  at (0, 1.5) {$m$};
  \node (m1) at (-1.5, 0) {$m_1$};
  \node (m2) at (1.5, 0) {$m_2$};
  \node (m3) at (0, -1.5) {$m_3$};

  \draw[->] (m) -- (m1);
  \draw[->] (m) -- (m2);
  \draw[->, dashed] (m1) -- (m3);
  \draw[->, dashed] (m2) -- (m3);
\end{tikzpicture}
\end{center}

Una relación \(R\) es \emph{confluente} si su clausura
reflexiva-transitiva \(R^*\) tiene la propiedad del diamante.

\begin{Shaded}
\begin{Highlighting}[]
\KeywordTok{def}\NormalTok{ confluent (R : relation X) : }\KeywordTok{Prop}\NormalTok{ :=}
\NormalTok{  diamond (clos\_refl\_trans R)}
\end{Highlighting}
\end{Shaded}

\textbf{Lema 1.2.1.} Si una relación \(R\) tiene la propiedad del
diamante, entonces su clausura transitiva \(R^+\) también la tiene.

El enunciado correspondiente en Lean es:

\begin{Shaded}
\begin{Highlighting}[]
\KeywordTok{lemma}\NormalTok{ diamond\_clos\_trans \{R : relation X\} :}
\NormalTok{    diamond R →}
\NormalTok{    diamond (clos\_trans R)}
\end{Highlighting}
\end{Shaded}

\bookmarksetup{startatroot}

\chapter{\texorpdfstring{Cálculo
\(\lambda\)}{Cálculo \textbackslash lambda}}\label{cuxe1lculo-lambda}

Para representar términos del cálculo \(\lambda\) utilizaremos la
representación \emph{locally nameless}. En ella, las variables ligadas
se representan mediante índices de de Bruijn, mientras que las variables
libres conservan un nombre. De esta manera, la sintaxis de los términos
está dada por

\[
t ::= 
\mathsf{bvar} \, i  \mid 
\mathsf{fvar} \, x  \mid 
\mathsf{abs}  \, t  \mid 
\mathsf{app}  \, t \, t  
\]

donde \(i \in \mathbb{N}\) es un índice de de Bruijn y \(x\) es el
nombre de una variable libre. En Lean tomamos los nombres como cadenas
de caracteres:

\begin{Shaded}
\begin{Highlighting}[]
\NormalTok{inductive trm : }\KeywordTok{Type}
\NormalTok{  | bvar : ℕ → trm}
\NormalTok{  | fvar : String → trm}
\NormalTok{  | abs  : trm → trm}
\NormalTok{  | app  : trm → trm → trm}
\end{Highlighting}
\end{Shaded}

Para simplificar la escritura, usamos cadenas de caracteres para
variables libres y naturales para variables ligadas:

\begin{Shaded}
\begin{Highlighting}[]
\NormalTok{instance : Coe String trm where}
\NormalTok{  coe := fvar}

\NormalTok{instance (n : ℕ) : OfNat trm n where}
\NormalTok{  ofNat := bvar n}
\end{Highlighting}
\end{Shaded}

Por ejemplo, el término \(\lambda x. \, x \, y\) puede escribirse en
Lean como \texttt{abs\ (app\ 0\ "y")}.

\subsection{Variables libres}\label{variables-libres}

El conjunto de variables libres de un término \(t\), denotado por
\(\operatorname{fv} (t)\), se define recursivamente mediante \[
\begin{array}{l c l}
\operatorname{fv}(\mathsf{bvar}\, i) & \equiv & \emptyset \\ 
\operatorname{fv}(\mathsf{fvar}\, x) & \equiv & \{x\} \\ 
\operatorname{fv}(\mathsf{app}\, t_1\, t_2) & \equiv & \operatorname{fv}(t_1) \cup \operatorname{fv}(t_2) \\ 
\operatorname{fv}(\mathsf{abs}\, t) & \equiv & \operatorname{fv}(t)
\end{array}
\] Su implementación en Lean es:

\begin{Shaded}
\begin{Highlighting}[]
\KeywordTok{def}\NormalTok{ fv (t : trm) : Finset String :=}
  \KeywordTok{match}\NormalTok{ t }\KeywordTok{with}
\NormalTok{  | bvar \_    =\textgreater{} ∅}
\NormalTok{  | fvar x    =\textgreater{} \{x\}}
\NormalTok{  | app t₁ t₂ =\textgreater{} fv t₁ ∪ fv t₂}
\NormalTok{  | abs u     =\textgreater{} fv u}
\end{Highlighting}
\end{Shaded}

\section{2.1. Apertura}\label{apertura}

En la representación \emph{locally nameless}, una abstracción tiene la
forma \(\mathsf{abs} \, t\), donde el cuerpo \(t\) utiliza índices de de
Bruijn para referirse a las variables ligadas. Para trabajar con este
cuerpo resulta útil reemplazar las referencias a la abstracción exterior
por una variable libre. A esta operación la llamaremos \emph{apertura}.

Para definirla, utilizamos una operación auxiliar
\(\{k\rightsquigarrow x\}\,t\), donde \(k\) indica el índice que se
reemplaza por la variable libre de nombre \(x\). Al entrar en una
abstracción, el parámetro \(k\) aumenta en uno. La operación se define
recursivamente por

\[
\begin{array}{l c l}
\{k \rightsquigarrow x\} \, (\mathsf{bvar}\, i) & \equiv & \text{si } (i = k) \text{ entonces } (\mathsf{fvar}\, x) \text{ si no } (\mathsf{bvar}\, i) \\ 
\{k \rightsquigarrow x\} \, (\mathsf{fvar}\, y) & \equiv & \mathsf{fvar}\, y \\ 
\{k \rightsquigarrow x\} \, (\mathsf{app}\, t_1 \, t_2) & \equiv & \mathsf{app}\, (\{k \rightsquigarrow x \} \, t_1) \, (\{k \rightsquigarrow x\} \, t_2) \\ 
\{k \rightsquigarrow x\} \, (\mathsf{abs}\, t) & \equiv & \mathsf{abs}\, (\{(k + 1) \rightsquigarrow x\} \, t)
\end{array}
\]

En Lean:

\begin{Shaded}
\begin{Highlighting}[]
\KeywordTok{def}\NormalTok{ open\_rec (k : ℕ) (x : String) (t : trm) : trm :=}
  \KeywordTok{match}\NormalTok{ t }\KeywordTok{with}
\NormalTok{  | bvar i    =\textgreater{} if i = k then fvar x else bvar i}
\NormalTok{  | fvar y    =\textgreater{} fvar y}
\NormalTok{  | app t₁ t₂ =\textgreater{} app (open\_rec k x t₁) (open\_rec k x t₂)}
\NormalTok{  | abs t     =\textgreater{} abs (open\_rec (k + 1) x t)}

\KeywordTok{notation} \StringTok{"\{"}\NormalTok{ k }\StringTok{" \textasciitilde{}\textgreater{} "}\NormalTok{ x }\StringTok{"\} "}\NormalTok{ t =\textgreater{} open\_rec k x t}
\end{Highlighting}
\end{Shaded}

La apertura del cuerpo \(t\) con una variable libre \(x\) se define
tomando \(k=0\): \[
t^x \equiv \{ 0 \rightsquigarrow x \} \, t
\]

En Lean:

\begin{Shaded}
\begin{Highlighting}[]
\KeywordTok{def}\NormalTok{ open\_var (t : trm) (x : String) : trm :=}
\NormalTok{  \{0 \textasciitilde{}\textgreater{} x\} t}

\NormalTok{infixl:80 }\StringTok{" \^{} "}\NormalTok{ =\textgreater{} open\_var}
\end{Highlighting}
\end{Shaded}

Así, \(t^x\) representa el resultado de aplicar la apertura al cuerpo de
una abstracción \(\mathsf{abs} \, t\), utilizando la variable libre
\(x\).

Por ejemplo, consideremos en Lean la abstracción \texttt{abs\ t}, con
\texttt{t\ =\ app\ (abs\ (app\ 0\ 1))\ 0}. Al abrir su cuerpo con el
nombre \texttt{"x"}, obtenemos:

\begin{Shaded}
\begin{Highlighting}[]
\KeywordTok{example}\NormalTok{ :}
\NormalTok{  (app (abs (app 0 1)) 0) \^{} }\StringTok{"x"}\NormalTok{ = (app (abs (app 0 }\StringTok{"x"}\NormalTok{)) }\StringTok{"x"}\NormalTok{) := }\KeywordTok{rfl}
\end{Highlighting}
\end{Shaded}

El índice \texttt{0} que aparece dentro de la abstracción interior
permanece sin cambios, pues hace referencia a dicha abstracción. En
cambio, el índice \texttt{1} que aparece dentro de ella y el índice
\texttt{0} exterior hacen referencia a la abstracción exterior y son
reemplazados por \texttt{"x"}.

\section{2.2. Clausura local}\label{clausura-local}

El tipo \texttt{trm} admite expresiones cuyos índices de de Bruijn no
hacen referencia a ninguna abstracción, como ocurre en \texttt{bvar\ 0}
y en \texttt{abs\ 1}. En el primer caso, el índice no se encuentra bajo
ninguna abstracción y, en el segundo, no corresponde a la única
abstracción presente, mientras que \texttt{abs\ 0} sí representa el
término identidad \(\lambda x.\,x\).

Diremos que un término es localmente cerrado si todos los índices de de
Bruijn que aparecen en él se encuentran ligados por alguna abstracción.
La proposición \(\mathsf{lc} \, t\) expresa que el término \(t\) es
localmente cerrado. Este predicado se define mediante las reglas

\[
\inferrule*[right=\textsc{lc-var}]
  { }
  {\mathsf{lc} \, (\mathsf{fvar}\,x)}
\qquad
\inferrule*[right=\textsc{lc-app}]
  {\mathsf{lc} \, t_1 \\ \mathsf{lc} \, t_2}
  {\mathsf{lc} \, (\mathsf{app} \, t_1 \, t_2)}
\qquad
\inferrule*[right=\textsc{lc-abs}]
  {\forall x \notin L, \quad \mathsf{lc} \, (t^x)}
  {\mathsf{lc} \, (\mathsf{abs}\,t)}
\]

donde \(L\) es algún conjunto finito de nombres.

En Lean, esta definición se expresa de la siguiente manera:

\begin{Shaded}
\begin{Highlighting}[]
\NormalTok{inductive lc : trm → }\KeywordTok{Prop}
\NormalTok{  | lc\_var (x : String) :}
\NormalTok{      lc (fvar x)}
\NormalTok{  | lc\_app \{t₁ t₂ : trm\} :}
\NormalTok{      lc t₁ → lc t₂ → lc (app t₁ t₂)}
\NormalTok{  | lc\_abs (L : Finset String) (t : trm) :}
\NormalTok{      (∀ x ∉ L, lc (t \^{} x)) →}
\NormalTok{      lc (abs t)}
\end{Highlighting}
\end{Shaded}

El constructor \texttt{lc\_var} establece que toda variable libre es
localmente cerrada. El constructor \texttt{lc\_app} indica que una
aplicación es localmente cerrada cuando lo son sus dos componentes. Por
último, \texttt{lc\_abs} emplea \emph{cuantificación cofinita}: exige
elegir un conjunto finito \texttt{L} y demostrar que \texttt{t\ \^{}\ x}
es localmente cerrado para todo nombre \texttt{x} que no pertenezca a
\texttt{L}. La definición no incluye un constructor para términos de la
forma \texttt{bvar\ i}, pues un índice aislado no está ligado por
ninguna abstracción y, por tanto, no es localmente cerrado.

\section{2.3. Cierre}\label{cierre}

Mientras que la apertura reemplaza determinadas referencias ligadas por
una variable libre, el cierre reemplaza las apariciones libres de un
nombre por índices de de Bruijn. En particular, dado un término \(t\) y
un nombre \(x\), permite construir el cuerpo de una abstracción que liga
esas apariciones de \(x\).

Para definirla, utilizamos una operación auxiliar
\(\{k \leftsquigarrow x \} \, t\), donde \(k\) indica el índice de de
Bruijn que debe asignarse a las variables libres con nombre \(x\). Se
define recursivamente por

\[
\begin{array}{l c l}
\{k \leftsquigarrow x\} \, (\mathsf{bvar}\, i) & \equiv & \mathsf{bvar}\, i \\ 
\{k \leftsquigarrow x\} \, (\mathsf{fvar}\, y) & \equiv & \text{si } (x = y) \text{ entonces } (\mathsf{bvar}\, k) \text{ si no } (\mathsf{fvar}\, y) \\ 
\{k \leftsquigarrow x\} \, (\mathsf{app}\, t_1 \, t_2) & \equiv & \mathsf{app}\, (\{k \leftsquigarrow x \} \, t_1) \, (\{k \leftsquigarrow x\} \, t_2) \\ 
\{k \leftsquigarrow x\} \, (\mathsf{abs}\, t) & \equiv & \mathsf{abs}\, (\{(k + 1) \leftsquigarrow x\} \, t)
\end{array}
\]

En Lean:

\begin{Shaded}
\begin{Highlighting}[]
\KeywordTok{def}\NormalTok{ close\_var\_rec (k : ℕ) (x : String) (t : trm) : trm :=}
  \KeywordTok{match}\NormalTok{ t }\KeywordTok{with}
\NormalTok{  | bvar i    =\textgreater{} bvar i}
\NormalTok{  | fvar y    =\textgreater{} if x = y then bvar k else fvar y}
\NormalTok{  | app t₁ t₂ =\textgreater{} app (close\_var\_rec k x t₁) (close\_var\_rec k x t₂)}
\NormalTok{  | abs t     =\textgreater{} abs (close\_var\_rec (k + 1) x t)}

\KeywordTok{notation} \StringTok{"\{"}\NormalTok{ k }\StringTok{" \textless{}\textasciitilde{} "}\NormalTok{ x }\StringTok{"\} "}\NormalTok{ t =\textgreater{} close\_var\_rec k x t}
\end{Highlighting}
\end{Shaded}

El cierre de un término \(t\) respecto de una variable libre \(x\) se
define tomando \(k=0\): \[
{}^{\textbackslash x} t \equiv \{0 \leftsquigarrow x \} \, t
\] En Lean:

\begin{Shaded}
\begin{Highlighting}[]
\KeywordTok{def}\NormalTok{ close\_var (t : trm) (x : String) : trm :=}
\NormalTok{  \{0 \textless{}\textasciitilde{} x\} t}

\NormalTok{infixl:80 }\StringTok{" \^{}* "}\NormalTok{ =\textgreater{} close\_var}
\end{Highlighting}
\end{Shaded}

Así, \({}^{\textbackslash x} t\) reemplaza las variables libres con
nombre \(x\) por índices de de Bruijn que, al formar
\(\mathsf{abs} \, ({}^{\textbackslash x} t)\), hacen referencia a esta
abstracción.

Por ejemplo:

\begin{Shaded}
\begin{Highlighting}[]
\KeywordTok{example}\NormalTok{ :}
\NormalTok{  (app (abs (app 0 }\StringTok{"x"}\NormalTok{)) }\StringTok{"x"}\NormalTok{) \^{}* }\StringTok{"x"}\NormalTok{ = (app (abs (app 0 1)) 0) := }\KeywordTok{rfl}
\end{Highlighting}
\end{Shaded}

En este ejemplo, la aparición libre de \texttt{"x"} exterior se
reemplaza por \texttt{0}, mientras que la situada dentro de la
abstracción se reemplaza por \texttt{1}.

\bookmarksetup{startatroot}

\chapter{\texorpdfstring{Reducción
\(\eta\)}{Reducción \textbackslash eta}}\label{reducciuxf3n-eta}

Una reducción es, sencillamente, una relación binaria sobre el conjunto
de todos los términos \(\lambda\). La única reducción que nosotros
trataremos es \(\eta\), cuya definición usual es la siguiente: \[
\inferrule*[right=\textsc{eta-red}]
  {x \notin \operatorname{fv}(t)}
  {\lambda x. \, t \, x
   \longrightarrow_{\eta} t}
\qquad
\inferrule*[right=\textsc{eta-app-1}]
  {t_1 \longrightarrow_{\eta} t_1'}
  {t_1 \, t_2
   \longrightarrow_{\eta}
   t_1' \, t_2}
\] \[
\inferrule*[right=\textsc{eta-app-2}]
  {t_2 \longrightarrow_{\eta} t_2'}
  {t_1 \, t_2
   \longrightarrow_{\eta}
   t_1 \, t_2'}
\qquad
\inferrule*[right=\textsc{eta-abs}]
  {t \longrightarrow_{\eta} t'}
  {\lambda x. t
   \longrightarrow_{\eta}
   \lambda x. t'}
\]

Las reglas \textsc{eta-abs}, \textsc{eta-app-1} y \textsc{eta-app-2},
que permiten reducir dentro del cuerpo de una abstracción o en
cualquiera de los componentes de una aplicación, son llamadas
\emph{reglas de congruencia}.

Llamaremos \emph{redex} \(\eta\) a un término de la forma
\(\lambda x.\,t\,x\) con \(x\notin\operatorname{fv}(t)\). Si \(t\) y
\(t'\) son términos tales que \(t \longrightarrow_{\eta} t'\), entonces
decimos que \(t\) se \emph{reduce} a \(t'\) en un solo paso.

Para las clausuras de la reducción \(\eta\), usaremos las siguientes
notaciones habituales:

\begin{enumerate}
\def\labelenumi{\arabic{enumi}.}
\item
  \(\longrightarrow_{\eta}^{=}\) denota la clausura reflexiva de
  \(\longrightarrow_\eta\). Así, \(t \longrightarrow_{\eta}^{=} t'\)
  significa que \(t=t'\) o que \(t \longrightarrow_\eta t'\).
\item
  \(\longrightarrow_{\eta}^{+}\) denota la clausura transitiva de
  \(\longrightarrow_\eta\). Escribimos
  \(t \longrightarrow_{\eta}^{+} t'\) cuando \(t\) se reduce a \(t'\) en
  uno o más pasos.
\item
  \(\twoheadrightarrow_{\eta}\) denota la clausura reflexiva-transitiva
  de \(\longrightarrow_\eta\). Escribimos
  \(t \twoheadrightarrow_\eta t'\) si \(t=t'\) o si una sucesión finita
  de pasos lleva de \(t\) a \(t'\).
\end{enumerate}

Como \(\eta\) es la única reducción que consideraremos, a partir de
ahora omitiremos el subíndice \(\eta\) en la notación de la reducción y
de sus clausuras.

Para formalizar la reducción \(\eta\) mediante la representación
\emph{locally nameless}, debemos adaptar las reglas anteriores.

En la regla \textsc{eta-red}, el término \(\lambda x.\,t\,x\) se
representa como \(\mathsf{abs}\,(\mathsf{app}\,t\,0)\), donde el índice
\(0\) hace referencia a la abstracción exterior. La condición de que
\(t\) sea localmente cerrado garantiza que todos sus índices hacen
referencia a abstracciones contenidas en el propio \(t\) y, en
particular, que ninguno hace referencia a la abstracción exterior que se
elimina. Así se expresa, en esta regla, la condición de que la variable
ligada no aparezca libre en \(t\).

Las reglas se formulan de modo que la reducción relacione únicamente
términos localmente cerrados. Al igual que la definición de
\(\mathsf{lc}\), la regla \textsc{eta-abs} emplea cuantificación
cofinita. Obtenemos así las siguientes reglas:

\[
\inferrule*[right=\textsc{eta-red}]
  {\mathsf{lc} \, t}
  {\mathsf{abs} \, (\mathsf{app} \, t \, 0)
   \longrightarrow t}
\qquad
\inferrule*[right=\textsc{eta-app-1}]
  {\mathsf{lc} \, t_2 \\ 
    t_1 \longrightarrow t_1'}
  {\mathsf{app}\,t_1\,t_2
   \longrightarrow
   \mathsf{app}\,t_1'\,t_2}
\]

\[
\inferrule*[right=\textsc{eta-app-2}]
  {\mathsf{lc} \, t_1 \\
   t_2 \longrightarrow t_2'}
  {\mathsf{app}\,t_1\,t_2
   \longrightarrow
   \mathsf{app} \, t_1 \, t_2'}
\qquad
\inferrule*[right=\textsc{eta-abs}]
  {\forall x\notin L, \quad
   t^x \longrightarrow t'^x}
  {\mathsf{abs}\,t
   \longrightarrow
   \mathsf{abs}\,t'}
\]

En Lean, esta relación se define inductivamente de la siguiente manera:

\begin{Shaded}
\begin{Highlighting}[]
\NormalTok{inductive eta : relation trm}
\NormalTok{  | eta\_red \{t : trm\} :}
\NormalTok{      lc t →}
\NormalTok{      eta (abs (app t (bvar 0))) t}
\NormalTok{  | eta\_app1 \{t₁ t₁\textquotesingle{} t₂ : trm\} :}
\NormalTok{      lc t₂ →}
\NormalTok{      eta t₁ t₁\textquotesingle{} →}
\NormalTok{      eta (app t₁ t₂) (app t₁\textquotesingle{} t₂)}
\NormalTok{  | eta\_app2 \{t₁ t₂ t₂\textquotesingle{} : trm\} :}
\NormalTok{      lc t₁ →}
\NormalTok{      eta t₂ t₂\textquotesingle{} →}
\NormalTok{      eta (app t₁ t₂) (app t₁ t₂\textquotesingle{})}
\NormalTok{  | eta\_abs (L : Finset String) \{t t\textquotesingle{} : trm\} :}
\NormalTok{      (∀ x, x ∉ L → eta (t \^{} x) (t\textquotesingle{} \^{} x)) →}
\NormalTok{      eta (abs t) (abs t\textquotesingle{})}
\end{Highlighting}
\end{Shaded}

\textbf{Lema 3.0.1.} Si \(t \longrightarrow t'\), entonces \(t'\) es
localmente cerrado.

En Lean:

\begin{Shaded}
\begin{Highlighting}[]
\KeywordTok{lemma}\NormalTok{ eta\_lc\_right \{t t\textquotesingle{} : trm\} :}
\NormalTok{    eta t t\textquotesingle{} → lc t\textquotesingle{}}
\end{Highlighting}
\end{Shaded}

\section{3.1. Reglas de congruencia para la clausura
reflexiva}\label{reglas-de-congruencia-para-la-clausura-reflexiva}

Las reglas de congruencia de la reducción \(\eta\) se extienden a su
clausura reflexiva. A continuación presentamos las dos que utilizaremos
en la demostración de la confluencia de la reducción \(\eta\).

\textbf{Lema 3.1.1.} \[
\inferrule*[right=\textsc{refl-app-1}]
  {\mathsf{lc} \, t_2 \\ 
    t_1 \longrightarrow^{=} t_1'}
  {\mathsf{app}\,t_1\,t_2
   \longrightarrow^{=}
   \mathsf{app}\,t_1'\,t_2}
\]

\begin{Shaded}
\begin{Highlighting}[]
\KeywordTok{lemma}\NormalTok{ refl\_app1 \{t₁ t₁\textquotesingle{} t₂ : trm\} :}
\NormalTok{    lc t₂ →}
\NormalTok{    clos\_refl eta t₁ t₁\textquotesingle{} →}
\NormalTok{    clos\_refl eta (app t₁ t₂) (app t₁\textquotesingle{} t₂)}
\end{Highlighting}
\end{Shaded}

\textbf{Lema 3.1.2.} \[
\inferrule*[right=\textsc{refl-app-2}]
  {\mathsf{lc} \, t_1 \\
   t_2 \longrightarrow^{=} t_2'}
  {\mathsf{app}\,t_1\,t_2
   \longrightarrow^{=}
   \mathsf{app} \, t_1 \, t_2'}
\]

\begin{Shaded}
\begin{Highlighting}[]
\KeywordTok{lemma}\NormalTok{ refl\_app2 \{t₁ t₂ t₂\textquotesingle{} : trm\} :}
\NormalTok{    lc t₁ →}
\NormalTok{    clos\_refl eta t₂ t₂\textquotesingle{} →}
\NormalTok{    clos\_refl eta (app t₁ t₂) (app t₁ t₂\textquotesingle{}) }
\end{Highlighting}
\end{Shaded}

\section{3.2. Análisis de las
reducciones}\label{anuxe1lisis-de-las-reducciones}

La forma del término inicial permite identificar las reglas con las que
puede obtenerse una reducción.

Si \(\mathsf{app}\,t_1\,t_2\longrightarrow s\), hay dos posibilidades:

\begin{enumerate}
\def\labelenumi{\arabic{enumi}.}
\item
  Se aplica la regla \textsc{eta-app-1}: \(t_2\) es localmente cerrado y
  existe un término \(t_1'\) tal que \(t_1\longrightarrow t_1'\) y
  \(s=\mathsf{app}\,t_1'\,t_2\).
\item
  Se aplica la regla \textsc{eta-app-2}: \(t_1\) es localmente cerrado y
  existe un término \(t_2'\) tal que \(t_2\longrightarrow t_2'\) y
  \(s=\mathsf{app}\,t_1\,t_2'\).
\end{enumerate}

Si \(\mathsf{abs}\,t\longrightarrow s\), las posibilidades son las
siguientes:

\begin{enumerate}
\def\labelenumi{\arabic{enumi}.}
\item
  Se aplica la regla \textsc{eta-red}: existe un término \(u\)
  localmente cerrado tal que \(t=\mathsf{app}\,u\,0\) y \(s=u\).
\item
  Se aplica la regla \textsc{eta-abs}: existen un conjunto finito de
  nombres \(L\) y un término \(t'\) tales que
  \(t^x\longrightarrow t'^x\) para todo \(x\notin L\) y
  \(s=\mathsf{abs}\,t'\).
\end{enumerate}

En particular, consideremos una reducción
\(\mathsf{abs}\,(\mathsf{app}\,t\,0)\longrightarrow s\), donde \(t\) es
localmente cerrado. Las dos posibilidades para esta reducción se
presentan de la siguiente manera.

Si se aplica \textsc{eta-red}, obtenemos directamente \(s=t\): \[
\inferrule*[right=\textsc{eta-red}]
  {\mathsf{lc}\,t}
  {\mathsf{abs}\,(\mathsf{app}\,t\,0)\longrightarrow t}.
\]

Si se aplica \textsc{eta-abs}, tenemos \(s=\mathsf{abs}\,u\) para algún
término \(u\), y la reducción se obtiene mediante \[
\inferrule*[right=\textsc{eta-abs}]
  {\forall x\notin L,\quad
   (\mathsf{app}\,t\,0)^x\longrightarrow u^x}
  {\mathsf{abs}\,(\mathsf{app}\,t\,0)
   \longrightarrow\mathsf{abs}\,u},
\] donde \(L\) es un conjunto finito de nombres.

En este último caso, el resultado es una abstracción
\(\mathsf{abs}\,u\). El siguiente lema permite determinar la forma de su
cuerpo a partir de la condición que exige la regla \textsc{eta-abs}.

\textbf{Lema 3.2.1.} Sea \(t\) un término localmente cerrado y sea \(L\)
un conjunto finito de nombres. Si
\((\mathsf{app}\,t\,0)^x\longrightarrow u^x\) para todo \(x\notin L\),
entonces existe un término \(t'\) tal que \(t\longrightarrow t'\) y
\(u=\mathsf{app}\,t'\,0\).

En Lean:

\begin{Shaded}
\begin{Highlighting}[]
\KeywordTok{lemma}\NormalTok{ eta\_redex\_body \{L : Finset String\} \{t u : trm\} :}
\NormalTok{    lc t →}
\NormalTok{    (∀ x ∉ L, eta ((app t (bvar 0)) \^{} x) (u \^{} x)) →}
\NormalTok{    ∃ t\textquotesingle{}, eta t t\textquotesingle{} ∧ u = app t\textquotesingle{} (bvar 0)}
\end{Highlighting}
\end{Shaded}





\end{document}
